\chapter{Discusión}

Nos enfocamos en el estudio de un caso de interacción abrupto (Caso 1) y otro de interacción suave (Caso 2), en los cuales se utiliza la misma CME precursora pero interactuando con diferentes CMEs sucesoras. Los resultados obtenidos permiten interpretar, más allá de la comparación cuantitativa presentada en la Sección \ref{análisisderesultados}, el papel que juegan los parámetros cinemáticos y la morfología de las CMEs en la viabilidad y naturaleza de su interacción, así como el alcance y las limitaciones del código desarrollado para representarla de forma realista. Sin embargo, no es posible caracterizar el tipo de interacción según el coeficiente de restitución $e$, debido a que no es posible diferenciar las velocidades de las CMEs después de su interacción.

La región más amplia de valores $(a_d, \tau_d)$ que permiten interacción en el Caso 1, frente a la región más reducida del Caso 2 (Figuras \ref{fig:advstd} y \ref{fig:advstd2} respectivamente), sugieren que una aceleración abrupta y de corta duración favorece la interacción incluso cuando la CME sucesora no llega a superar en velocidad a la precursora ($V_{\text{CME-2}} > 0,86\,V_{\text{CME-1}}$). Mientras que una aceleración suave y prolongada exige que la sucesora efectivamente iguale o supere dicha velocidad ($V_{\text{CME-3}} > V_{\text{CME-1}}$). Por lo tanto, una CME que crece abruptamente en aceleración tiende a estabilizar su velocidad rápidamente, de modo que el umbral de interacción depende de alcanzar rápidamente una velocidad de alcance por parte de la CME sucesora, modulada por los parámetros $a_d$ y $\tau_d$. 

La diferencia notable en los coeficientes de correlación entre el tiempo y altura de interacción con los parámetros $a_d$ y $\tau_d$ en el Caso 1 ($R^2 \sim 0,89$–$1,0$) frente al Caso 2 ($R^2 \sim 0,017$–$0,089$) está relacionada con la fase de aceleración principal de la CME sucesora. Cuando el tiempo de ascenso $\tau_r$ de la sucesora es largo, el punto de encuentro entre CMEs deja de estar gobernado por su fase de decaimiento y pasa a depender casi enteramente de la cinemática de la precursora y del retraso entre erupciones. Por lo tanto, la fase de propagación de la CME sucesora no influye en el comportamiento de la interacción, lo que implica que para CMEs sucesoras de ascenso lento predecir el punto de interacción a partir únicamente de sus parámetros $a_d$ y $\tau_d$ podría ser poco informativo, ya que los parámetros dominantes serían $\tau_r$, $a_r$ y su tiempo de retraso.

La interacción mostrada en las Figuras \ref{cinematicapolarconjunta} y \ref{cinematicapolarconjunta2} reproduce cualitativamente rasgos tales como un frente de baja densidad seguido de una región de alta densidad (Figura \ref{Serietiempocaso1}), junto con la región de menor densidad en comparación con la de fondo tras el paso de las CMEs debido al arrastre de material sobre el medio. Asimismo, el hecho de que el contacto entre frentes y colas ocurra tempranamente (10 horas en el Caso 1 y 18 h en el Caso 2, en $<30\,R_\odot$) mientras que la interacción de los centros se desarrolla mucho después (26 y 70 h), reproduce razonablemente bien la escala temporal reportada para eventos reales de interacción CME-CME \parencite[e.g.,][]{mishra-2014}. 

Los resultados sugieren que la cinemática de dos fases (Ecuación \eqref{ecuaceleración}) combinada con la expansión autosimilar \parencite[e.g.,][]{savani-2011, riley-2004} es suficiente para capturar la geometría temprana de la interacción, aunque sin procesos magnéticos como la reconexión\index{reconexión magnética}. Por esto, la estructura resultante de la interacción obedece al comportamiento de la presión dinámica, siendo una densidad de energía cinética del flujo, la cual se incrementa en las regiones donde la velocidad relativa entre las CMEs es alta, generando allí un fuerte gradiente de presión y, en consecuencia, una región de compresión con mayor densidad; mayor en el Caso 1 ($\sim 10^3$ protones cm$^{-3}$) en relación al Caso 2 ($\sim 3\times 10^2$ protones cm$^{-3}$). Sin embargo, en la región trasera de la CME sucesora ocurre el proceso inverso, la disminución de la velocidad reduce la presión dinámica local dando lugar a una región de expansión y de menor densidad, como ocurre en el Caso 1, ver Figura \ref{Serietiempocaso1}.

En ambos casos es evidente que la interacción incrementa la densidad de la CME resultante, y presumiblemente el campo magnético, en relación a las CMEs individuales. Esto es consistente con la estadística mencionada en la \autoref{CMESEPs}, según la cual las CMEs gemelas producen eventos de \acp{SEP}\index{Partículas energéticas Solares} grandes con mucha mayor frecuencia que las CMEs simples ($\sim 60\%$ frente a $\sim 20\%$), probablemente debido a un alto nivel de complejidad magnética presente en la región de interacción. Lo que refuerza la idea de que la interacción de CMEs es un factor que determina en mayor medida el comportamiento del clima espacial, en relación a CMEs simples o individuales.

Es necesario precisar los alcances y las restricciones físicas del código: \textit{i}) El campo de densidad y velocidad de la CME resultante se construye mediante un promedio ponderado de densidades y velocidades, no mediante la solución de las ecuaciones de conservación de momento y energía. Por lo que, magnitudes como el campo magnético amplificado en las regiones de interacción, la formación de un choque inverso o la reconexión interna no emergen del modelo, sino que solo pueden inferirse indirectamente a partir del comportamiento de la densidad y la velocidad simuladas. \textit{ii}) El \ac{SW}\index{Solar Wind} de fondo se mantiene homogéneo y estacionario, lo cual simplifica la heliosfera\index{Heliosfera} interna real, donde la estructura de flujos rápidos y lentos puede modificar sustancialmente la probabilidad de interacción \parencite[e.g.,][]{aschwanden-2019}. Finalmente, \textit{iii}) la morfología autosimilar impuesta a partir de \textcite[e.g.,][]{savani-2011} no permite el aplanamiento ni las deformaciones no autosimilares reportadas en interacciones reales \parencite[e.g.,][]{lugaz-2005,xiong-2006}, aunque sí introduce asimetría angular mediante las perturbaciones de Fourier y el ruido angular incorporado en los perfiles.

Estas limitaciones no invalidan los resultados cualitativos obtenidos, pero limitan su interpretación. El modelo es apropiado para explorar las condiciones cinemáticas bajo las cuales dos CMEs homólogas\index{homologas, Coronal Mass Ejection} pueden interactuar y la deformación de sus perfiles de densidad y velocidad al hacerlo, mientras que preguntas sobre la intensificación del campo magnético, la formación de choques inversos o la producción cuantitativa de \acp{SEP}\index{Partículas energéticas Solares} requieren un tratamiento magnetohidrodinámico en el código.